\section{A Geometric Reference Baseline}
\label{sec:method}

The ceiling motivates a change of scope. A crack is a mark at a fixed location
on a rigid, near-planar surface that carries texture the crack does not; we
identify a crack by \emph{where it is on its wall} once the two photographs are
aligned. We release this as a deliberately simple \emph{reference baseline} for
\textsc{CrackID}, with no learned component beyond the shared crack segmenter.

\subsection{Alignment driven by non-crack texture}
A wall is planar, so the homography $x' \sim Hx$ mapping photograph $I_a$ into
$I_b$'s frame is the correct model. The design decision that matters is
\textbf{what drives the fit}: we dilate the crack masks by $15$\,px and exclude
those pixels from detection, so $H$ comes from wall texture alone. This is not an
optimisation but what makes the projected position an independent identity cue ---
were cracks to participate, a long crack could dominate the correspondence set
and pull the fit toward the wrong crack. Detection runs on copies downscaled to
$1600$\,px under CLAHE equalisation; we use SIFT (up to $8000$ features), a
$0.75$ ratio test, and MAGSAC at $4$\,px requiring $25$ inliers.

\subsection{A cascade for two surface families}
One front end does not serve both surfaces, and the failure is categorical: on
the smooth-plaster walls, crack-masked SIFT succeeds on \emph{exactly zero}
pairs. We try three in order: (i)~\textbf{SIFT on masked wall texture}, for
stippled render; (ii)~\textbf{SIFT with cracks restored}, triggered only when
masking starves detection below $300$ keypoints (on plaster the crack is the only
feature in frame), accepting the circularity stage~(i) avoids, and recording
which front end produced each result; (iii)~\textbf{ECC}, direct photometric
alignment over a $320 \to 640 \to 1280$\,px pyramid, on the large-scale shading
that survives on flat paint. Every candidate $H$ is sanity-checked (reject
mirrored/folded quadrilaterals, area scale outside $[1/16,16]$, collapsed sides);
its inlier RMS sets downstream tolerances.

\subsection{Projection and Chamfer agreement}
Crack instances are connected components; pixels from $I_a$ are projected by $H$
and compared with instances of $I_b$. Chamfer distance, not IoU, is the
criterion: cracks are a few pixels wide, so a $3$\,px registration error drives
IoU to zero for a correctly aligned pair. The reported score is the
\textbf{symmetric mean Chamfer distance} in pixels
\begin{equation}
D(a,b) = \tfrac{1}{2}\!\left[
\frac{1}{|a|}\sum_{p \in a} d(p, b) \;+\;
\frac{1}{|b|}\sum_{q \in b} d(q, a)\right],
\label{eq:chamfer}
\end{equation}
over the union bounding box, with cheap rejection when boxes are far apart. We
also report a scale-free alternative --- \textbf{Chamfer coverage}, the symmetric
fraction of pixels within $\tau$ --- because Eq.~(\ref{eq:chamfer}) is a per-pair
pixel distance whose scale differs per pair, so a single global threshold makes
ranking partly measure calibration (Sec.~\ref{sec:chamferablation}).

\subsection{Explicit non-answers, and the licensed negative}
Correspondence is a one-to-one assignment per image pair with an explicit
rejection threshold, not a per-query argmax --- two cracks must not both claim
the same crack; unmatched instances are reported as \emph{disappeared} or
\emph{newly appeared}. The Hungarian solve is part of \emph{evaluation}, applied
identically to every method's score matrix. Specific to this formulation is the
evidence a \textbf{licensed negative} can carry: with $I_a$ projected into
$I_b$'s frame, the system observes a region visible and empty of the recorded
crack --- positive evidence of absence, categorically stronger than the low
similarity a crop embedding produces. A projected crack falling outside the
gallery photograph is a confident rejection, counted separately from unscorable
pairs; this is the mechanism behind the open-set gap
(Sec.~\ref{sec:coverage}).